% Original analytic teaching diagram, 2026-11-12.
% Probability tables on {0,1,2} as the triangle Delta_2 with vertices delta_0=(0,0),
% delta_1=(60,0), delta_2=(30,51.96) (mm). A table (p0,p1,p2) sits at p0*delta_0+p1*delta_1+p2*delta_2.
% Mean-one constraint p0=p2: segment from delta_1 (0,1,0) to (1/2,0,1/2) = midpoint (15,25.98) of edge delta_0 delta_2.
% Uniform table (1/3,1/3,1/3) = centroid (30,17.32). Edge p2=0 is delta_0--delta_1.
\definecolor{VB}{HTML}{7188A3}\definecolor{VBF}{HTML}{E5EAF0}
\definecolor{VG}{HTML}{829C8B}\definecolor{VGF}{HTML}{E8EEE9}
\definecolor{VY}{HTML}{B4A16F}\definecolor{VYF}{HTML}{F1ECDD}
\definecolor{VR}{HTML}{A7787D}\definecolor{VRF}{HTML}{F0E5E6}
\definecolor{VP}{HTML}{9588A5}\definecolor{VPF}{HTML}{ECE8F0}
\definecolor{VI}{HTML}{3B4252}\definecolor{VM}{HTML}{6B7280}\definecolor{VL}{HTML}{4C4D4F}
\begin{tikzpicture}[x=1mm,y=1mm,>=Stealth,font=\figurefont\fontsize{8}{10}\selectfont,
lab/.style={text=VM,font=\figurefont\fontsize{7.5}{9}\selectfont}]
\coordinate (D0) at (0,0);
\coordinate (D1) at (60,0);
\coordinate (D2) at (30,51.96);
\fill[VBF] (D0) -- (D1) -- (D2) -- cycle;
\draw[VB,line width=.6pt] (D0) -- (D2) -- (D1);
\draw[VL,line width=1.6pt] (D0) -- (D1);
\draw[VR,line width=1.4pt] (D1) -- (15,25.98);
\fill[VR] (30,17.32) circle (1.1);
\foreach \p in {(D0),(D1),(D2)} {\fill[VI] \p circle (0.8);}
\node[text=VI,anchor=north east] at (D0) {$\delta_0$};
\node[text=VI,anchor=north west] at (D1) {$\delta_1$};
\node[text=VI,anchor=south] at (D2) {$\delta_2$};
\node[lab,anchor=north] at (30,-1) {$p_2=0$};
\node[text=VR,anchor=south east] at (14.5,26.5) {$(\tfrac12,0,\tfrac12)$};
\node[text=VR,anchor=west] at (31,19.8) {均匀表};
\node[text=VR,anchor=south west] at (22.5,24.2) {$\mathbb{E}_P[X]=1$};
\end{tikzpicture}
